\knowledge{1.1.1}
{交集要求同时属于两集合，并集要求至少属于一个，差集$A\setminus B$要求属于$A$且不属于$B$。补集依赖指定全集。区间端点用开闭括号标明；De Morgan律为$(A\cup B)^c=A^c\cap B^c$及$(A\cap B)^c=A^c\cup B^c$。}
{能把绝对值不等式化为区间，准确计算交、并、差、补及对称差。能用任取元素的双向包含证明集合相等，特别检查端点是否取到。}

\knowledge{1.1.2}
{上界不必属于集合；上确界是最小上界。非空有上界实数集存在上确界。$s=\sup A$等价于：所有$a\in A$满足$a\le s$，且任意$\varepsilon>0$均有$a\in A$使$s-\varepsilon<a$。下确界有对称表述。}
{能同时验证“是界”与“可任意逼近”，区分确界、最大元和最小元。能按数列奇偶项或单调性求界，并用定义处理集合平移、并集的确界。}

\knowledge{1.1.3}
{函数由定义域、陪域和对应规则确定，值域包含在陪域内。单射表示不同输入有不同输出，满射表示每个陪域元素都有原像，双射同时满足两者。集合原像$f^{-1}(B)=\{x:f(x)\in B\}$不要求反函数存在。}
{判断两个函数是否相同前先比较定义域。能由$f(x_1)=f(x_2)$检验单射，解$f(x)=y$检验满射，并求指定集合的原像。}

\knowledge{1.1.4}
{实线性映射同时满足$f(x+y)=f(x)+f(y)$与$f(\lambda x)=\lambda f(x)$，其中$\lambda$为任意实数。对$\mathbb R\to\mathbb R$的映射，这等价于$f(x)=ax$。仿射函数$ax+b$只有在$b=0$时为线性映射。}
{能按两条定义逐项验证，并推导$f(0)=0$、$f(-x)=-f(x)$。能用$f(x)=xf(1)$确定映射，区分保持仿射组合与实线性。}

\knowledge{1.1.5}
{复合$(f\circ g)(x)=f(g(x))$，定义域为$\{x\in D_g:g(x)\in D_f\}$。复合一般不满足交换律。根式要求被开方数非负，对数要求真数为正，分母不能为零。}
{先写内层的定义域，再加入外层对内层值的限制。能分别求$f\circ g$与$g\circ f$并比较；化简后仍保留原来的定义域限制。}

\knowledge{1.1.6}
{双射存在反函数，反函数的定义域是原函数的陪域，值域是原定义域。实区间上的严格单调函数对其值域可逆。反函数图像与原图像关于$y=x$对称；$f^{-1}(x)$与$\frac{1}{f(x)}$含义不同。}
{能先限制定义域以获得单射，再解方程、交换变量并写出反函数的定义域和值域。能验证两种复合均为相应集合上的恒等映射。}

\knowledge{1.1.7}
{双曲函数定义为$\sinh x=\frac{e^x-e^{-x}}{2}$、$\cosh x=\frac{e^x+e^{-x}}{2}$，满足$\cosh^2x-\sinh^2x=1$。$\sinh$为奇函数，$\cosh$为偶函数，$\tanh x=\frac{\sinh x}{\cosh x}$。}
{能求复合初等函数的实定义域，用指数定义证明双曲恒等式。求反双曲函数时先检查单调分支、值域和对数真数条件。}

\knowledge{1.2.1}
{$\lim_{n\to\infty}a_n=a$表示：任意$\varepsilon>0$，存在正整数$N$，使所有$n>N$都有$|a_n-a|<\varepsilon$。$N$可依赖$\varepsilon$，不能依赖正在考查的$n$。有限项改变不影响极限。}
{能从$|a_n-a|<\varepsilon$倒推出可用的$N$，再按量词顺序完整证明。能用不同极限的子列或不趋零的误差证明不收敛。}

\knowledge{1.2.2}
{收敛数列极限唯一且有界；极限四则运算中，商的分母极限须非零。最终有$a_n\le b_n$且两列收敛，则$a\le b$；严格不等式取极限通常只保留非严格形式。有界性不能单独推出收敛。}
{能检查四则运算条件，利用有界性与趋零因子控制乘积。能给出“正项极限为零”“有界但不收敛”等反例，避免把必要条件当充分条件。}

\knowledge{1.2.3}
{单调有界数列收敛。夹逼准则要求上下两列趋于同一极限。Cauchy准则为：任意$\varepsilon>0$，存在$N$使$m,n>N$时$|a_m-a_n|<\varepsilon$，无需预先知道极限。}
{递推数列先证明定义有效、单调和有界，再令极限满足递推关系并选取符合范围的根。能用尾项估计证明Cauchy性，用比较估计证明部分和有界。}

\knowledge{1.3.1}
{$\lim_{x\to a}f(x)=L$要求$0<|x-a|<\delta$时$|f(x)-L|<\varepsilon$，不要求$f(a)$有定义。双侧有限极限存在当且仅当左右极限存在且相等。无穷极限要用任意大的阈值描述。}
{能按$\varepsilon$--$\delta$定义证明简单极限，分别计算左右极限。能构造两列不同函数值极限的自变量证明振荡函数无极限，并写出无穷极限的量词。}

\knowledge{1.3.2}
{有限极限满足和、差、积运算，商的分母极限须非零。只在穿孔邻域中相等的函数仍可有相同极限。$\lim_{x\to a}f(x)=0$等价于$\lim_{x\to a}|f(x)|=0$；乘积趋零不能单独推出指定因子趋零。}
{能约去非极限点处的公因子后求极限，并说明约分的适用范围。能从和差反求两个函数的极限，为不可逆的推论构造反例。}

\knowledge{1.3.3}
{基本极限为$\lim_{u\to0}\frac{\sin u}{u}=1$和$\lim_{u\to0}(1+u)^{\frac1u}=e$（取$1+u>0$）。由此得$\lim_{u\to0}\frac{\ln(1+u)}{u}=1$。三角极限中的角度用弧度。}
{能把内层变量完整换成$u$并补齐倍数。指数型极限先检查底数为正，必要时取对数，将指数与对数的乘积化为基本极限。}

\knowledge{1.3.4}
{夹逼准则可控制有界振荡因子。Heine定理要求对任意趋于目标点、避开目标点且在定义域内的数列，函数值均趋于同一极限。单调且有相应界的函数在区间端点有有限单侧极限。}
{能找到统一的绝对值上界并夹逼趋零，能用两列检验否定极限存在。证明单调函数端点极限时能用确界的逼近性质，并区分有限极限与无穷极限。}

\knowledge{1.4.1}
{同一趋近过程中，有限个无穷小之和、无穷小与有界函数的乘积均趋零。两个无穷小的乘积也是无穷小。项数随变量增长的和需要另外估计，不能直接套用有限和结论。}
{能先指出哪一因子有界、哪一因子趋零，再给出估计。能识别结论成立所需的条件，并通过除以比较量验证小$o$关系。}

\knowledge{1.4.2}
{在同一趋近过程$x\to a$中，且比较量最终非零时，$\alpha=o(\beta)$表示$\lim_{x\to a}\frac{\alpha(x)}{\beta(x)}=0$，$\alpha\sim\beta$表示$\lim_{x\to a}\frac{\alpha(x)}{\beta(x)}=1$。同阶通常指比值极限为非零有限常数，等价是其中常数等于1的情形。}
{能用比值极限判断阶，掌握$\sin u\sim u$、$1-\cos u\sim\frac{u^2}{2}$、$\ln(1+u)\sim u$等基本关系，并随内层变量同步替换。}

\knowledge{1.4.3}
{等价无穷小可在乘除结构中代换，前提是比较关系成立且分母最终非零。加减结构可能消去首项，直接分别代换可能丢失真正的主导项；此时可用有理化、恒等变形或更高精度展开。}
{能先识别运算结构，写出每一步等价关系再求比值极限。遇到相减先检查首项是否抵消，能够解释误差项为何影响最终阶数。}

\knowledge{1.4.4}
{$\lim_{x\to a}f(x)=+\infty$表示在该趋近过程中最终超过任意给定正阈值；无界只表示不存在统一有限界。$\lim_{x\to a}|f(x)|=+\infty$允许符号变化，但不允许反复回到有限范围。最终为正时，$\lim_{x\to a}f(x)=+\infty$等价于$\lim_{x\to a}\frac{1}{f(x)}=0$。}
{能分别描述左右趋近时的符号与大小，利用主导项或阈值定义证明无穷极限。能用保留有界子列的反例区分无界、绝对值趋于无穷和趋于正无穷。}

\knowledge{1.5.1}
{在定义域的聚点$x_0$连续要求$\lim_{x\to x_0}f(x)=f(x_0)$。有限双侧极限存在但点值不匹配为可去间断；有限左右极限不等为跳跃间断；振荡或无穷行为形成第二类间断。}
{能先算左右极限再分类，判断修改一个点值能否补成连续。含有理、无理分段时，能分别沿两类点逼近并找出全部连续点。}

\knowledge{1.5.2}
{连续性对和、差、积保持，对商要求分母不为零；复合连续要求内层值落在外层定义域。分段连接点连续要求左右极限与点值相等。连续函数在非零点附近保持符号。}
{能逐层检查定义域，用拼接条件确定参数。能以$\max(f,g)=\frac{f+g+|f-g|}{2}$等恒等式证明连续性，用取$\varepsilon<|f(x_0)|$证明局部保号。}

\knowledge{1.5.3}
{闭有界区间上的连续函数有界并取得最大、最小值，且具有介值性。端点异号可保证内部零点，但唯一性需额外论证。处处为正时，由最小值定理可得到严格正的统一下界。}
{能逐项检查连续、闭区间和端点条件。能把存在性问题改写为辅助函数有零点，用单调性证明唯一性，并说明区间不闭时结论可能失效。}

\knowledge{1.5.4}
{一致连续要求：每个$\varepsilon>0$都有一个对全域通用的$\delta>0$，使$|x-y|<\delta$时$|f(x)-f(y)|<\varepsilon$。闭区间连续函数一致连续；Lipschitz估计也是充分条件。}
{能给出不依赖位置的差值估计。证明不一致连续时，构造两列使$\lim_{n\to\infty}|x_n-y_n|=0$而函数值差不趋零，避免只凭导数无界直接下结论。}

\knowledge{1.5.5}
{闭区间上的自映射若满足$|T(x)-T(y)|\le q|x-y|$且$0\le q<1$，则存在唯一不动点，迭代收敛。完备性、自映射和严格压缩均需检查；有不动点本身不保证任意迭代收敛。}
{能验证区间不被映出、确定压缩常数并求不动点。能利用几何衰减估计迭代误差与次数，解释开区间或$q=1$时定理不能直接使用。本小节按授课进度选做。}

\knowledge{2.1.1}
{导数$f'(a)=\lim_{h\to0}\frac{f(a+h)-f(a)}{h}$，存在时要求左右差商趋于同一有限值。点处的特别定义必须放回差商，不能直接使用只对$x\ne a$成立的导函数表达式。}
{能从定义计算多项式、根式和分段函数的点导数。振荡因子先估计差商，再判断其是否趋零；能清楚区分函数连续与差商收敛。}

\knowledge{2.1.2}
{可导曲线在$(a,f(a))$的切线为$y-f(a)=f'(a)(x-a)$。$f'(a)\ne0$时法线斜率为$-\frac{1}{f'(a)}$；切线水平时法线竖直。平行切线条件由导数等于指定斜率得到。}
{能先求导、求接触点，再写切线与法线方程。能处理水平切线及法线竖直的情况，确认所求点属于原函数定义域。}

\knowledge{2.1.3}
{可导必连续，连续不一定可导；函数可导也不保证导函数连续。分段点可导须先连续，再使左右导数相等。导函数在该点连续还需要邻近导数的极限等于该点导数。}
{能依次检查连续、可导和导函数连续，分别给出判据与反例。证明可导推出连续时，利用函数增量等于差商与自变量增量的乘积。}

\knowledge{2.1.4}
{平均变化率为$\frac{\Delta y}{\Delta t}$，瞬时变化率为$\frac{dy}{dt}$。速度是位移对时间的一阶导数，加速度是二阶导数；边际量为总量对相应自变量的导数。导数的单位是因变量单位除以自变量单位。}
{能识别自变量、函数与单位，在指定时刻代入导数。收入等复合量先建立完整函数再求导，解释变化率的正负与实际含义。}

\knowledge{2.2.1}
{$(u\pm v)'=u'\pm v'$，$(uv)'=u'v+uv'$，$\left(\frac uv\right)'=\frac{u'v-uv'}{v^2}$，商法则要求$v\ne0$。证明乘积法则时要对增量拆项，并使用可导函数连续。}
{能准确识别各因子并展开法则，检查分母与对数等定义域。能从乘积法则推出倒数法则，也能用差商完成基本法则证明。}

\knowledge{2.2.2}
{链式法则为$(f(g(x)))'=f'(g(x))g'(x)$，每一层都要满足可导条件。正函数可先取对数，把乘积、商或变指数转成便于求导的形式；对数求导后需乘回原函数。}
{能按外层到内层逐层求导且不漏内层因子。能用给定点值计算复合函数导数，对$x^x$一类变指数函数使用对数求导。}

\knowledge{2.2.3}
{设$y=f(x)$在所取分支可逆，且$f'(x_0)\ne0$，在通常的反函数可导条件下，$(f^{-1})'(f(x_0))=\frac{1}{f'(x_0)}$。分母导数为零时不能直接用此公式判断有限导数。}
{能先找原像点$x_0$再代入倒数公式，并注明反三角函数所用分支。能从反函数关系推导对数和反三角导数，单独检查公式失效的特殊点。}

\knowledge{2.2.4}
{常用公式包括$(e^x)'=e^x$、$(\ln x)'=\frac1x$（$x>0$）、$(\sinh x)'=\cosh x$、$(\cosh x)'=\sinh x$。根式、对数及反函数表达式求导前要确认实定义域与可导范围。}
{能结合积法则和链式法则计算初等复合函数导数。能从指数定义推导双曲导数，并在化简根式时保留正负号和区间条件。}

\knowledge{2.2.5}
{高阶导数由逐次求导定义。Leibniz公式为$(uv)^{(n)}=\sum_{k=0}^n\binom nk u^{(k)}v^{(n-k)}$。指数函数导数有幂因子，三角导数有四步周期，多项式足够高阶导数为零。}
{能从前几阶发现模式并用归纳法证明，正确处理阶乘、符号与$n$的范围。对多项式乘指数能用Leibniz公式压缩非零项。}

\knowledge{2.2.6}
{对$F(x,y)=0$，若$F$足够光滑且$F_y\ne0$，可局部将$y$视为$x$的函数，$y'=-\frac{F_x}{F_y}$。隐式求导过程中每个$y$都要按复合函数处理，二阶导数需再次使用链式法则。}
{能先验证点在曲线上，再计算一阶、二阶导数。能说明局部可解条件，检查除以$y$或其他表达式时是否遗漏特殊点。}

\knowledge{2.2.7}
{若$x=x(t)$、$y=y(t)$且$x'(t)\ne0$，则$\frac{dy}{dx}=\frac{y'}{x'}$，$\frac{d^2y}{dx^2}=\frac{y''x'-y'x''}{(x')^3}$。二阶求导是再对$x$求导，需要再除以$x'(t)$。}
{能把指定参数代入导数及坐标，写出切线方程。能在计算前检查$x'(t)$是否为零，并解释二阶公式的推导及适用条件。}

\knowledge{2.2.8}
{相关变化率先把各变化量视为同一时间$t$的函数，建立几何或物理约束，再两边对$t$求导。长度、面积、体积和压力等变化率单位不同，方向由正负号体现。}
{能列清已知瞬时量和待求导数，先求导再代入某时刻的数据。能由相似关系、勾股关系、体积公式或守恒条件建立正确方程，并解释下降速度的符号。}

\knowledge{2.3.1}
{一元函数在一点可微等价于可导。增量可写为$\Delta y=f'(x)\Delta x+o(\Delta x)$，微分$dy=f'(x)dx$是线性主部。余项与$\Delta x$之比趋零，线性系数唯一。}
{能区分实际增量$\Delta y$与微分$dy$，在指定点和$dx$下计算微分。能通过余项比值验证可微性，说明绝对值函数在尖点不能用同一个线性主部表示。}

\knowledge{2.3.2}
{微分满足和差、积商法则；例如$d(uv)=u\,dv+v\,du$。一阶形式不变性指即使$x$为$t$的函数，仍有$dy=f'(x)dx$，此时$dx=x'(t)dt$。商的分母仍须非零。}
{能把同一微分分别用$du$、$dx$或$dt$表示，保持自变量层次清楚。能由链式法则解释一阶形式不变性，并检查各函数的可导条件。}

\knowledge{2.3.3}
{当$x$为独立变量时，$d^ny=f^{(n)}(x)(dx)^n$。当$x$依赖另一独立变量时，$d^2y=f''(x)(dx)^2+f'(x)d^2x$；不能普遍省去第二项，高阶微分通常不具有一阶的形式不变性。}
{每次计算先注明独立变量，能分别按$x$和$t$求高阶微分。能通过链式求导写出额外项，并用具体复合函数检验两种计算方式相符。}

\knowledge{2.3.4}
{当$h$较小时，$f(a+h)\approx f(a)+f'(a)h$，一般误差为$o(h)$。相对误差用$\frac{\Delta y}{y}$或其线性近似$\frac{dy}{y}$描述。若要给出严格误差上界，需要额外的余项估计。}
{能选择数值容易计算且靠近目标的基点，写出函数、增量和近似结果。能由幂函数微分估计相对误差，并明确近似等号与严格等号的区别。}

\knowledge{2.4.1}
{Fermat定理：内点处可导并取得局部极值，则导数为零。驻点只是候选点，未必有极值；不可导点及区间端点也可能取得极值，但不满足该定理的全部前提。}
{能列出驻点、不可导点和端点，随后另作极值判别。能从左右差商的符号证明Fermat定理，并用实例解释每个条件的必要性。}

\knowledge{2.4.2}
{Rolle定理要求闭区间连续、内部可导且端点等值。Lagrange定理给出$f(b)-f(a)=f'(\xi)(b-a)$。Cauchy定理给出$[f(b)-f(a)]g'(\xi)=[g(b)-g(a)]f'(\xi)$；改写比值时另查分母非零。}
{能完整核对区间条件并写出存在点所在范围。能构造辅助函数处理带$f'+f$的结论，用导数界估计函数差，避免把存在点当作任意点。}

\knowledge{2.4.3}
{L'Hospital法则适用于满足相应可导及分母导数非零等条件的$\frac00$或$\frac\infty\infty$型比值，并要求导数比具有相应极限。其他未定式需先变形，幂型通常先取对数。导数比无极限不能反推原比无极限。}
{每次使用前先判定未定式并检查条件，连续使用时逐次重查。能把$0\cdot\infty$转为比值，把变指数转为对数极限，并识别本法则不适用的情形。}

\knowledge{2.5.1}
{Taylor多项式系数为$\frac{f^{(k)}(a)}{k!}$。Peano余项形式为$f(a+h)=\sum_{k=0}^n\frac{f^{(k)}(a)h^k}{k!}+o(h^n)$；在相应高阶可导条件下，Lagrange余项为$\frac{f^{(n+1)}(\xi)h^{n+1}}{(n+1)!}$。}
{能区分局部阶数结论与可估计的误差界，写清展开点、阶数和条件。能说明$o(h^n)$未必是$O(h^{n+1})$，利用低阶导数消失确定首个可能的非零项。}

\knowledge{2.5.2}
{零点展开包括$e^x=1+x+\frac{x^2}{2}+\cdots$，$\sin x=x-\frac{x^3}{6}+\cdots$，$\cos x=1-\frac{x^2}{2}+\cdots$，$\ln(1+x)=x-\frac{x^2}{2}+\cdots$。$(1+x)^\alpha$使用下降连乘系数，$\arctan x$只含交替奇次项。}
{能写出指定有限阶的系数与Peano余项，牢记奇偶性和阶乘。能在内层变量趋零时进行代入，并把余项换算到最终变量对应的阶数。}

\knowledge{2.5.3}
{极限计算应展开到分子消去后首个非零项，余项精度须足以与分母比较。数值近似要用Lagrange余项或其他严格界控制误差；Taylor多项式的最高非零次数不一定等于用于估计的定理阶数。}
{能根据分母阶数决定展开深度，对复合表达式保留必要项。能在指定区间估计高阶导数，明确近似值及绝对误差上界，避免只给形式展开。}

\knowledge{2.6.1}
{区间上可导且$f'\ge0$可推出单调不减，$f'>0$可推出严格递增。严格递增函数在某些点导数仍可为零。导数符号判断需按定义域与临界点分区间进行。}
{能求导、排列临界点并给出符号表及单调区间。能用中值定理证明单调性，处理导数偶尔为零时的严格单调问题，并构造逆命题反例。}

\knowledge{2.6.2}
{局部极值可用一阶导数左右变号判别。若$f'(x_0)=0$且$f''(x_0)>0$则为严格局部极小，$f''(x_0)<0$则为严格局部极大；$f''(x_0)=0$时该判据不作结论。不可导点也需检查。}
{能根据导数符号变化判断极大与极小，并写出点和函数值。能在二阶判据失效时回到定义或一阶判别，用二阶展开证明充分条件。}

\knowledge{2.6.3}
{闭区间连续函数的最值可在驻点、不可导点与端点之间比较。开区间或无界区间还需研究边界极限并确认最值是否取得。应用题的可行域来自正值、长度和几何约束。}
{能先将约束化为一元函数和正确的定义域，再列全候选点并比较。能区分最大值与上确界，解释所得参数是否满足实际条件。}

\knowledge{2.6.4}
{在光滑情形，$f''>0$表示图像向上弯曲，$f''<0$表示向下弯曲。拐点要求曲线在该点两侧改变弯曲方向；$f''=0$或二阶导数不存在仅给出候选点，还需检查连续性与两侧符号。}
{能按定义域求出全部候选点，逐个验证二阶导数的符号变化并给出坐标。能说明二阶导数为零不足以保证拐点，必要时计算拐点切线。}

\knowledge{2.7.1}
{平面曲线曲率描述单位弧长上的切线转角变化，$\kappa=\left|\frac{d\varphi}{ds}\right|$。曲率非负，其单位为长度的倒数。$\kappa>0$时曲率半径$R=\frac1\kappa$；曲率为零时没有有限曲率半径。}
{能从转角与弧长关系直接求曲率，解释直线、圆及恒曲率的几何含义。能在曲率和曲率半径之间换算并核对单位。}

\knowledge{2.7.2}
{对$y=f(x)$，$\kappa=\frac{|f''|}{[1+(f')^2]^{\frac32}}$。对正则参数曲线，$\kappa=\frac{|x'y''-y'x''|}{[(x')^2+(y')^2]^{\frac32}}$，要求速度向量非零。曲率半径只在曲率非零时取倒数。}
{能分别计算一、二阶导数后代入正确公式，保留绝对值并检查分母。能把参数点转换为曲线位置，说明曲率为零或参数不正则时的特殊情况。}

\knowledge{3.1.1}
{累积量可由“局部密度或变化率乘微小区间长度”求和并取极限得到。面积由高度累积，质量由线密度累积，位移由速度累积，净流入量由流入率减流出率累积；路程使用速率。}
{能说明积分变量、上下限、被积函数及单位，先正确列式。能写出等分矩形和，并用上下和或已知求和公式理解从离散和到积分的过程。}

\knowledge{3.1.2}
{Riemann积分是分割细度趋零时，任取采样点的和$\sum f(\xi_i)\Delta x_i$的共同极限。对有界函数，上下和之差可任意小是可积判据。只考查某一种采样方式不足以一般地证明可积。}
{能从等分和识别$\Delta x$、采样点和区间，利用求和公式算出定义积分。能用上下和分析阶梯函数，并借有理、无理数稠密性判断典型不可积函数。}

\knowledge{3.1.3}
{普通闭区间Riemann可积函数必须有界；连续函数可积，有界且只有有限个间断点的函数也可积。有界本身不保证可积。有限个点值的修改不改变有界可积函数的积分。}
{能先判定有界性，再使用适当的充分条件。能区分普通积分与反常积分，用上下和或小区间覆盖解释有限异常点对积分的影响。}

\knowledge{3.1.4}
{定积分满足线性、区间可加与保序。$|\int_a^bf|\le\int_a^b|f|$（$a<b$）；连续函数满足积分中值定理$\int_a^bf=(b-a)f(c)$。对称区间上奇函数积分为零，偶函数积分为半区间的两倍。}
{能利用对称性和上下界估计积分，写清可积或连续等条件。能证明严格正的连续函数在正长度区间积分为正，并结合介值性说明积分平均值能够取到。}

\knowledge{3.2.1}
{在通常的连续被积函数与原函数条件下，Newton--Leibniz公式为$\int_a^bf(x)dx=F(b)-F(a)$。使用前需确认整个积分区间上成立，跨越奇点不能直接代端点。同一连通区间上的两个原函数只差常数。}
{能求原函数并正确代入上下限，按绝对值或分段点拆开积分。能检查区间内奇点，利用中值定理证明原函数差为常数。}

\knowledge{3.2.2}
{若$f$连续，则$F(x)=\int_a^xf(t)dt$满足$F'(x)=f(x)$。复合上下限求导为$\frac d{dx}\int_{u(x)}^{v(x)}f(t)dt=f(v(x))v'(x)-f(u(x))u'(x)$，需满足相应连续、可导条件。}
{能区分积分哑变量与求导变量，计算双变限和复合上限导数。能用差商与连续性证明基本定理，并检查积分函数在特殊点的高阶可导性。}

\knowledge{3.2.3}
{不定积分表示同一连通区间上的全体原函数$F+C$。常用表包括幂函数、$\frac1x$、指数、三角和反三角积分；其中$\int\frac1x\,dx=\ln|x|+C$只在不跨零的区间上使用。}
{能逆用求导公式并写出积分常数，通过求导核对结果。能说明对数、根式及反三角原函数的适用区间，区分不同连通区间上可独立选择的常数。}

\knowledge{3.3.1}
{第一类换元利用$\int f(g(x))g'(x)dx=F(g(x))+C$。第二类换元以$x=\varphi(t)$化简根式等结构。定积分换元时变量、微分与上下限必须同时改变，并检查换元的定义域和区间条件。}
{能识别内层及其导数，选择直接换元或三角换元。能同步计算新上下限，处理根式中的绝对值，并通过还原变量或求导验证结果。}

\knowledge{3.3.2}
{分部积分公式为$\int u\,dv=uv-\int v\,du$；定积分还需端点项。常见选择让对数、反三角或多项式求导后简化；循环积分要联立解出，递推积分要写初值与适用阶数。}
{能明确列出$u,dv,du,v$并处理端点。能求解循环出现的原积分，建立递推关系并检查首个适用的$n$，不遗漏积分常数。}

\knowledge{3.3.3}
{有理函数可先分解分母作部分分式，二次式可配方，含根式可通过换元有理化，三角恒等式可改变积分结构。连续函数在区间上有原函数，但原函数未必能写成有限的初等函数组合。}
{能根据结构选择方法并核对定义域，通过求导验证化简后的结果。能区分“存在原函数”“可用初等式表示”和“定积分可定义”这几件事。}

\knowledge{3.4.1}
{微元法先把对象分成薄片或短段，写$dQ$再累加。常见结构为$dm=\rho(x)dx$、$dV=A(x)dx$、$dW=F(x)dx$；积分变量必须与截面、密度或距离的表达相匹配。}
{能说明微元的几何或物理意义，正确确定累加范围。能从微元单位检查总量单位，辨认截面积、半径、线密度和变化率，避免直接套用不匹配公式。}

\knowledge{3.4.2}
{平面面积为$\int|f-g|dx$，绕轴旋转用圆盘或圆环截面；图像弧长为$\int\sqrt{1+(y')^2}dx$。极坐标面积为$\frac12\int r^2d\theta$，需选择能正确覆盖区域的角度范围。}
{能先确定交点与积分区间，判断上下曲线及内外半径。能区分面积与有符号积分，选择合适的弧长或极坐标公式并检查是否重复覆盖。}

\knowledge{3.4.3}
{沿力方向的功由$F\,ds$累积。静水压力元为水密度乘$g$、深度和面积元；抽水功还需提升距离。质心$\bar x=\frac{\int x\,dm}{M}$，转动惯量$I=\int r^2dm$，其中$r$是到转轴的距离。}
{能选坐标并把深度、提升高度、力和密度写成该坐标的函数。能建立质量微元与力矩，核对功、合力、质心及惯量的不同单位与物理含义。}

\knowledge{3.5.1}
{无穷区间积分按截断极限定义，如$\int_a^{\infty}f=\lim_{R\to\infty}\int_a^Rf$。全实轴积分须在某有限点拆成两侧，各自收敛才算收敛；对称截断的主值不能替代普通反常积分。}
{能先写截断积分再取极限，判断所得值是否有限。能辨认振荡导致截断极限不存在的情形，并对双侧无穷区间分别检查。}

\knowledge{3.5.2}
{被积函数在有限端点无界时，使用避开奇点的单侧截断极限。若区间内部有奇点，须从两侧分别积分并要求各自收敛；若两端均有奇点，也要分别检查。}
{能找全奇点并正确分段，写清趋近方向，再计算原函数端点极限。能识别可积的无界奇性，避免用一次端点代入掩盖发散。}

\knowledge{3.5.3}
{对$\int_1^{\infty}x^{-p}dx$，收敛条件为$p>1$。非负函数可用比较或正的有限极限比值判别；绝对收敛推出收敛。对数因子常通过$t=\ln x$转换为新的幂型尾积分。}
{能确认比较区间及函数非负性，判断无穷远的主导阶。能选择可积上界或发散下界，分别给出“收敛”与“绝对收敛”的结论及依据。}

\knowledge{3.5.4}
{对$\int_0^1x^{-p}dx$，收敛条件为$p<1$，与无穷远幂型的条件不同。有限奇点主要由距离奇点的幂次决定；对数因子可通过$t=-\ln x$转为无穷区间问题。}
{能在每个奇点附近比较主导量，选取正确的局部幂型模型。能处理三角函数的局部等价关系与对数修正，逐个检查端点后再判断整段积分。}

\knowledge{3.5.5}
{对实参数$s$，Gamma积分$\Gamma(s)=\int_0^{\infty}t^{s-1}e^{-t}dt$在$s>0$收敛。递推为$\Gamma(s+1)=s\Gamma(s)$，$\Gamma(1)=1$，$\Gamma(\frac12)=\sqrt\pi$；缩放由$u=\beta x$得到。}
{能分别检查零点和无穷远的收敛性，分部积分时说明边界项消失。能用递推计算正整数、半整数值，并在积分换元前列出参数允许范围。}

\knowledge{4.1.1}
{微分方程的阶数由最高阶导数决定。线性方程中未知函数及各阶导数只以一次线性组合出现，系数只依赖自变量；齐次性看独立项是否为零。初值用于选择解，解还需注明定义区间。}
{能识别阶数、线性与齐次，代入导数验证解。能区分含任意常数的通解与给定初值的特解，处理隐式解的函数分支和不可解点。第4章按授课进度选做。}

\knowledge{4.1.2}
{对$y'=a(x)b(y)$，在$b(y)\ne0$的区间可写$\frac{dy}{b(y)}=a(x)dx$积分。被除掉的$b(y)=0$可能给出平衡解，须单独检查。初值问题的唯一性需要相应条件，不能只凭方程形式保证。}
{能先列平衡解，再求非平衡解并代入初值。能写包含初始点的有效解区间，检查分母、对数和根式限制；能用实际解验证非唯一情形。}

\knowledge{4.1.3}
{标准一阶线性方程为$y'+p(x)y=q(x)$。在系数连续的区间取积分因子$\mu=e^{\int p(x)dx}$，便有$(\mu y)'=\mu q$，从而$y=\mu^{-1}(\int\mu q\,dx+C)$。}
{能先将最高阶导数系数化为1，注明除法后的区间。能计算积分因子、求出一个通解并代入初值，最后用原方程核验。}

\knowledge{4.1.4}
{Bernoulli方程$y'+py=qy^n$在适当非零解区间可令$z=y^{1-n}$化为线性；齐次型$y'=F(\frac yx)$可令$y=vx$（$x\ne0$）。含$x+y$或$e^y$等结构时，可用对应组合为新变量。}
{能说明换元的动机，正确计算新变量导数并解转换后的方程。换回后检查零解、被除掉的解和实数定义域，代入原方程检验。}

\knowledge{4.1.5}
{缺少$y$的二阶方程可令$p=y'$并以$x$为自变量；缺少$x$时，可在适当区间把$p$看成$y$的函数，用$y''=p\frac{dp}{dy}$降阶。直接给出最高阶导数时可逐次积分。}
{能按缺失变量选择代换，分清$\frac{dp}{dx}$与$\frac{dp}{dy}$。除以$p$或$y$前检查常值解和零点，积分后保留足够的独立常数并注明解区间。}

\knowledge{4.1.6}
{增长模型、冷却模型和线性阻力模型由变化率规律建立，混合模型遵守“存量变化率等于流入量率减流出量率”。Logistic模型含有限容量约束。初值与单位是模型的一部分。}
{能把文字转成方程和初值，区分浓度、总量、流率及体积。能求解并解释长期极限、平衡状态、翻倍时间或终端速度，检查结果在物理范围内。}

\knowledge{4.2.1}
{弹簧振动模型为$my''+cy'+ky=F(t)$，位移从平衡位置量起。串联RLC电路电荷方程为$Lq''+Rq'+\frac qC=E(t)$。含$y^2$、$(y')^2$等项的方程一般不为线性。}
{能对应质量、阻尼、恢复力和外力，确定符号及单位。能识别常系数或变系数、齐次或非齐次，并由物理关系建立高阶线性方程。}

\knowledge{4.2.2}
{齐次线性方程的解可作常系数线性组合。标准$n$阶线性方程在连续系数区间有$n$个线性无关的基本解，非齐次通解为齐次通解加一个特解。同一非齐次方程两解之差满足对应齐次方程。}
{能验证解、计算Wronskian，并在相应线性方程的条件下判定基本解组。能由一个特解与齐次基本解写出通解，解释各任意常数的作用。}

\knowledge{4.2.3}
{常系数齐次方程通过特征多项式求解。实根$r$的重数为$m$时对应$e^{rx},xe^{rx},\ldots,x^{m-1}e^{rx}$；共轭根$\alpha\pm i\beta$给出$e^{\alpha x}\cos\beta x$与$e^{\alpha x}\sin\beta x$，重根时再乘相应幂次。}
{能列出特征方程并按根型写全独立解，确认任意常数个数与阶数相符。能代入初值确定常数，用代回方程检查符号和指数。}

\knowledge{4.2.4}
{非齐次通解为$y_h+y_p$。对多项式、指数与三角组合的右端可用待定系数；若试探结构与齐次解重合，按对应特征根重数乘$x^m$修正。多个右端可分别求特解再相加。}
{能先求齐次解再选择完整试探式，检查共振与重数，代入确定系数。能区分特解中的待定系数和通解中的自由常数，并验证整个方程。}

\knowledge{4.2.5}
{Euler方程的系数按导数阶数配有$x$的幂。在$x>0$可试$y=x^m$或令$t=\ln x$，此时$xy'=Y_t$、$x^2y''=Y_{tt}-Y_t$。已知一个非零解时可令$y=y_1v$降阶。}
{能准确转换二阶导数并保留定义区间，处理实根、重根和复根。已知解降阶时能选取其非零区间，构造第二个独立解并核验。}

\knowledge{4.3.1}
{一阶线性系统为$X'=A(t)X+b(t)$，齐次系统有$b=0$。$n$维系统的初值指定$n$个分量。$n$阶标量方程可通过引入$y,y',\ldots,y^{(n-1)}$改写为$n$维一阶系统。}
{能写出矩阵、未知向量与非齐次项，识别分量间的非线性项。能逐分量验证向量解，并在标量方程与一阶系统之间转换。本节按星号内容选做。}

\knowledge{4.3.2}
{基本矩阵$\Phi$的列为线性无关的齐次向量解，满足$\Phi'=A\Phi$且可逆。齐次通解为$X=\Phi C$，初值解为$\Phi(t)\Phi(t_0)^{-1}X_0$。非齐次两解之差仍为齐次解。}
{能检查每列是否为解及行列式是否非零，写出常向量形式的通解。能代入初值并验证，使用线性性证明叠加关系且不把非齐次解任意叠加。}

\knowledge{4.3.3}
{常系数齐次系统$X'=AX$的特征对$(\lambda,v)$给出解$e^{\lambda t}v$。独立特征向量足够时可组合成通解；复特征对可取实、虚部，缺少特征向量的重根结构会产生$t e^{\lambda t}$等项。}
{能处理对角、对称、旋转及简单重根系统，保留所需独立常数。能代入初值，依据特征值实部或显式解讨论长期行为，并逐分量检验。}

\knowledge{4.3.4}
{非齐次线性系统通解为齐次通解加一个特解。设$\Phi$为齐次基本矩阵，令$X=\Phi c(t)$可得$c'=\Phi^{-1}b$，从而$X=\Phi(C+\int\Phi^{-1}b\,dt)$。与齐次项重合的驱动可能产生共振项。}
{能选择逐分量求解、消元或常数变易，并计算矩阵乘法的正确顺序。能由初值定常数，检验右端与初值同时满足，识别共振导致的$t$因子。}

\knowledge{4.3.5}
{耦合系统各分量遵守各自的流入减流出关系。总量导数由分量导数相加得到；总量守恒时可约束平衡状态。平衡解满足$AX+b=0$，零特征值常与一族平衡或守恒关系相关。}
{能由实际交换速率和浓度建立系统，写出初始存量。能求解、验证守恒量并讨论长期状态；存在平衡族时，用给定总量选择相应平衡，保持单位与非负范围一致。}
